\chapter{Foundations}\label{ch:foundations}

\section{Options, Black--Scholes and implied volatility}
Let $S$ denote the underlying, $F=F(t,T)$ its forward for delivery at $T$, and $D=D(t,T)$ the discount factor. Under the Black (forward) model, the price of a European call of strike $K$ and time-to-maturity $\tau=T-t$ is
\begin{equation}\label{eq:black}
  C(K,\tau;\sigma) \;=\; D\left[\,F\,N(d_+) - K\,N(d_-)\,\right],
  \qquad d_{\pm} \;=\; \frac{\log(F/K)}{\sigma\sqrt{\tau}} \pm \frac{\sigma\sqrt{\tau}}{2},
\end{equation}
with $N$ the standard normal cdf. Since $\partial C/\partial\sigma>0$ (the vega), \cref{eq:black} can be inverted: the \emph{implied volatility} $\iv(K,\tau)$ of a quote $C^{\mathrm{mkt}}$ is the unique root of $C(K,\tau;\sigma)=C^{\mathrm{mkt}}$. The empirical fact that $\iv$ varies with $(K,\tau)$---the \emph{smile} and the term structure---contradicts the constant-$\sigma$ assumption and gives rise to the implied volatility surface.

\section{Coordinates: log-moneyness and total variance}
Two changes of variable make the no-arbitrage theory tractable. The \emph{log-moneyness} $k=\log(K/F)$ centres the smile on the forward. The \emph{total variance}
\begin{equation}
  w(k,\tau) \;=\; \iv^2(k,\tau)\,\tau
\end{equation}
is the natural vertical coordinate: both static no-arbitrage conditions take their simplest form in $w$.

\section{Static no-arbitrage}\label{sec:noarb}
Following \cite{gatheral2014arbitrage, roper2010arbitrage}, a surface is free of static arbitrage if and only if each maturity slice is free of butterfly arbitrage and the family of slices is free of calendar arbitrage.

\begin{definition}[Butterfly arbitrage and the Durrleman function]
For a slice $w(\cdot)$ at fixed maturity, define
\begin{equation}\label{eq:durrleman}
  g(k) \;=\; \Big(1-\frac{k\,w'(k)}{2\,w(k)}\Big)^{2} \;-\; \frac{w'(k)^{2}}{4}\Big(\frac{1}{w(k)}+\frac14\Big) \;+\; \frac{w''(k)}{2}.
\end{equation}
The slice is free of butterfly arbitrage if and only if $g(k)\ge 0$ for all $k$ and call prices tend to zero at extreme strikes (for SVI, the wing-slope condition $b(1+|\rho|)<2$).
\end{definition}

The financial content of \cref{eq:durrleman} is the Breeden--Litzenberger identity: the risk-neutral density is proportional to $\partial^2 C/\partial K^2$, and in total-variance coordinates
\begin{equation}
  p(k) \;=\; \frac{g(k)}{\sqrt{2\pi\,w(k)}}\,\exp\!\Big(-\tfrac12 d_-(k)^2\Big),
  \qquad d_-(k) = -\frac{k}{\sqrt{w(k)}}-\frac{\sqrt{w(k)}}{2},
\end{equation}
so $g<0$ is exactly a negative implied density.

\begin{definition}[Calendar arbitrage]
The surface is free of calendar arbitrage if and only if $\tau\mapsto w(k,\tau)$ is non-decreasing for every $k$; equivalently, total-variance slices never cross.
\end{definition}

\Cref{eq:durrleman} involves only $w$, $w'$ and $w''$. This is the property exploited throughout the thesis: for parametric models the derivatives are closed-form; for neural models they are obtained \emph{exactly} by automatic differentiation, turning the no-arbitrage conditions into differentiable training penalties.

% FIGURE (export from NB02, Section 2): Vogt smile three panels
\begin{figure}[t]
  \centering
  \maybegraphics[width=\textwidth]{figures/nb02_vogt_arbitrage.png}
  \caption[The Vogt smile: hidden butterfly arbitrage]{The Vogt SVI parameters of \cite[Ex.~3.1]{gatheral2014arbitrage}. The total-variance slice (left) looks innocent, yet the Durrleman function dips below zero (centre) and the implied risk-neutral density turns negative (right): a genuine butterfly arbitrage invisible to the naked eye. Figure produced by our replication (Notebook~02).}
  \label{fig:vogt}
\end{figure}

\section{The reconstruction problem}
Given a day's quotes $\{(k_i,\tau_i,{\iv}_i)\}_{i=1}^{p}$, the problem studied in this thesis is to construct $\widehat{w}:\R\times(0,\infty)\to(0,\infty)$ that (i) fits the quotes, (ii) satisfies the conditions of \cref{sec:noarb}, and (iii) generalises to quotes not used in the construction. Conditions (i)--(iii) conflict; quantifying the trade-off across method families is the object of the thesis.
